% \section*{Выводы}
% \addcontentsline{toc}{section}{Выводы}

В ходе выполнения работы были получены следующие основные результаты.

\begin{enumerate}
    \item Разработан и формально описан фреймворк \textbf{Orbitals}~---
          целеориентированная RFM-сегментация пользователей на
          поведенческие архетипы (орбитали), обладающие двумя
          ключевыми свойствами: предсказательной способностью
          относительно ограниченной пожизненной ценности (lLTV) и
          чувствительностью к экспериментальным воздействиям.

    \item Предложен орбитальный оцениватель долгосрочного эффекта
          A/B-эксперимента на основе метрик $r_V$ и
          $r_{V\text{-corr}}$, формализующих долгосрочный эффект через
          перераспределение пользователей между орбиталями и их
          относительные ценности $\{V_c\}_{c \in \mathcal{C}}$.

    \item Сформулировано и эмпирически подтверждено мультипликативное
          допущение $V_c(t) = V_c \cdot V_o(t)$, позволяющее
          математически отделить сезонные колебания метрик от
          поведенческих сдвигов, вызванных продуктовым воздействием.
          На историческом окне в 1.5 года относительный порядок
          орбиталей по lLTV и масштаб различий между ними оставались
          устойчивыми.

    \item На корпусе из 21 значимого A/B-эксперимента крупной платформы
          электронной коммерции показано, что орбитальный оцениватель:
          \begin{itemize}
              \item корректно восстанавливает направление
                    долгосрочного эффекта~--- ни в одном из
                    рассмотренных экспериментов статистически значимый
                    положительный прогноз не соответствовал
                    статистически значимому отрицательному фактическому
                    эффекту (при уровне значимости $\alpha = 0.1$);
              \item снижает дисперсию оценки в $1.5$--$2.5$ раза по
                    сравнению с прямым наблюдением реализованной lLTV
                    и с сильной индивидуальной предсказательной моделью;
              \item обеспечивает повышенную чувствительность~--- в
                    ${\sim}21\%$ экспериментов обнаруживает значимый
                    эффект там, где краткосрочная метрика
                    \textit{Average GMV per User} его не различает.
          \end{itemize}

    \item На примере эксперимента BNPL продемонстрирована
          интерпретируемость метода: долгосрочный эффект объясняется
          конкретными изменениями в матрицах переходов~--- ростом
          вероятностей перехода в орбитали \texttt{active\&paid}
          ($+2.7\%$) и \texttt{whale} ($+8.7\%$) и снижением доли
          пользователей в низкоактивных состояниях ($-2.6\%$).

    \item Проведено сопоставление с двумя эталонными методами и
          эмпирически воспроизведён компромисс смещение--дисперсия:
          небольшое модельное смещение обменивается на существенное
          снижение дисперсии за счёт перехода в низкоразмерное
          пространство орбиталей.

    \item Реализована открытая программная библиотека и опубликован
          обезличенный набор данных более чем 1~млн пользователей,
          обеспечивающие воспроизводимость результатов.

    \item Установлены границы применимости метода: требование
          стабильности орбиталей во времени и условие симметричного
          воздействия внешних факторов на группы. Метод неприменим
          без дополнительной диагностики для интервенций, целенаправленно
          воздействующих на отдельные орбитали.
\end{enumerate}
